Low Earth orbit (LEO) has become significantly more congested over recent decades. This is down to a large increase in the number of satellites launched into LEO, but also the detection of smaller fragments. As a result, traditional conjunction assessment has fallen behind computationally. Many organisations are looking for machine learning (ML) to step in. This paper investigates the effectiveness of an ML model as a first stage filter for conjunction screening. Trained on real-world historical conjunction data, the distinction that this paper explores is the model’s ability to distinguish between satellite pairs that should generate a conjunction data message (CDM) and those pairs that occupy broadly similar orbital regions, despite not generating a CDM when propagated. To achieve this, 2 datasets were produced (containing only a satellite pairs’ orbital features). The positive dataset contained satellite pairs that had historically flagged a CDM. The negative dataset contained satellite pairs that had similar orbital features but did not generate a CDM within 7-days. The ML model, which was built on a random forest framework, was then evaluated on one key metric: recall rate. The recall rate of the ML model was determined by the number of incorrectly predicted ‘safe’ satellites – chosen specifically in order to operate a safe-heavy system, as a false positive was far less dangerous than a false negative in our context. When faced with a generic set of just under 300 satellite pairs, the model’s recall rate was 99.65\% resulting in 1 false negative encounter. When the dataset was then altered to match stricter non-conjunction criteria, the model’s recall rate fell to 97.80\% with 6 false negative predictions. Whilst this shows promise that an ML model can intervene as a first-stage filter in the conjunction assessment framework, it also suggests that further development of neural networks in the satellite safety context is required. Understandably so as 1 false negative prediction could cost a satellite operator millions of dollars. 